\chapter{When to Use Agents and When Not To}

Agents are powerful, but they are not the right abstraction for every problem. Many tasks are better solved by deterministic code, retrieval, a classifier, a workflow engine, or a single model call. The best agent architect knows when not to build an agent.

\section{Use an agent when the task is adaptive}

Agents are useful when the path to completion cannot be fully specified in advance. This happens when:

\begin{itemize}[leftmargin=*]
  \item the task is multi-step;
  \item the next step depends on observations;
  \item multiple tools may be needed;
  \item information is incomplete at the start;
  \item the environment changes;
  \item there are recoverable failures;
  \item the system must decide when to ask a human.
\end{itemize}

Examples include debugging a codebase, researching a topic, resolving a support case, investigating an incident, reconciling documents, or planning a trip with changing constraints.

\section{Do not use an agent when a simpler system works}

Do not use an agent when:

\begin{itemize}[leftmargin=*]
  \item a deterministic script is enough;
  \item a single API call answers the request;
  \item a classifier can route the case;
  \item a retrieval system can return the document directly;
  \item the workflow is fixed and does not need model decisions;
  \item ultra-low latency is required;
  \item errors are expensive and hard to detect;
  \item permissions cannot be safely scoped;
  \item success cannot be evaluated.
\end{itemize}

A model should not be used to add numbers when a calculator exists. An agent should not be used to update records when a deterministic workflow can do so safely.

\section{A decision ladder}

Use the simplest abstraction that meets the requirement.

\begin{enumerate}[leftmargin=*]
  \item Deterministic code.
  \item Search or database query.
  \item Classifier or extractor.
  \item Single LLM call.
  \item RAG plus single LLM call.
  \item Tool-calling assistant.
  \item Graph-based agent.
  \item Multi-agent system.
\end{enumerate}

Move up the ladder only when the lower level fails because the task requires adaptation, judgment, or multi-step interaction.

\section{Agent readiness checklist}

Before building an agent, ask:

\begin{itemize}[leftmargin=*]
  \item What is the success criterion?
  \item Which tools are needed?
  \item Which actions have side effects?
  \item Can the agent verify progress?
  \item What is the maximum budget?
  \item What should trigger human approval?
  \item What is the failure taxonomy?
  \item What evaluation dataset will be used?
  \item What is the rollback plan?
\end{itemize}

If these questions cannot be answered, the agent is not production-ready.

\section{Common overuse patterns}

\begin{description}[leftmargin=*]
  \item[Agent as expensive router.] A frontier model classifies tasks that a cheap classifier could route.
  \item[Agent as unreliable API wrapper.] The model chooses parameters for a fixed API call when a form would be safer.
  \item[Agent as search replacement.] The agent summarizes without reliable retrieval or citations.
  \item[Agent as policy engine.] The model decides authorization instead of infrastructure.
  \item[Agent as workflow engine.] The model improvises a fixed business process that should be encoded deterministically.
\end{description}

\section{Summary}

Agents are best for uncertain, multi-step, tool-rich workflows where observations determine the next step. They are not a replacement for deterministic software. The right design uses agents only where model judgment adds value and surrounds them with tools, knowledge, security, evaluation, and cost controls.
